\section{Finite-cover arithmetic without rotational symmetry}
\label{sec:rotation-free-arithmetic}
The circular proof of measurable phase rigidity used rotations to
eliminate a constant collision phase. The constant can instead be
included in the integer cocycle. This gives an arithmetic argument
which is unchanged by a noncircular deformation of the scatterer.

Let $T$ be an invertible ergodic collision map on a periodic planar
billiard, $\kappa:M\to\Z^2$ its lattice increment, and $\tau$ its flight
time. We use the regular local stable--unstable product rectangles
of the finite-horizon dispersing theory, with absolute continuity of
both conditional measures, as in \cite[Section 8]{Young}. The action
identity is \eqref{eq:v35-action-general}. Assume that on every finite
lattice cover $p\Lambda$, $p$ prime, the collision map is mixing.
These are geometric and ergodic hypotheses; no twisted spectrum or
local limit theorem is assumed.

\begin{lemma}[A finite character separates a proper subgroup]
\label{lem:v35-finite-character}
If $L$ is a proper subgroup of $\Z^d$, there are a prime $p$ and a
nonzero homomorphism $\chi:\Z^d\to\Z/p\Z$ which vanishes on $L$.
\end{lemma}
\begin{proof}
The quotient $\Z^d/L$ is a nonzero finitely generated abelian group.
If it has a nonzero free summand, reduction of one such coordinate
modulo any prime gives a quotient of order $p$. Otherwise one of its
nonzero finite cyclic summands has order divisible by a prime $p$;
reduction of that summand gives the required quotient. Compose with
the quotient map from $\Z^d$. This also covers subgroups of deficient
rank without assuming they have finite index.
\end{proof}

\begin{theorem}[Complete physical phase rigidity]
\label{thm:v35-rotation-free-rigidity}
Suppose a measurable circle-valued function $q$ satisfies
\begin{equation}\label{eq:v35-general-phase}
 q\circ T=e^{i(u\cdot\kappa+s+b\tau)}q,
       \qquad u\in\R^2,\ s,b\in\R.
\end{equation}
Then $b=0$, $s\in2\pi\Z$, $u\in(2\pi\Z)^2$, and $q$ is almost
everywhere constant. No rotational or reflection symmetry of the table
is required.
\end{theorem}
\begin{proof}
We first recall the measurable holonomy argument, including the step
which removes the roof frequency. Choose a positive-measure compact
regular product rectangle $P$ away from grazing. Conditional pair
measures on a common local stable or unstable leaf have both marginals
dominated by $C_P\nu$. Approximate $q$ in $L^1(\nu)$ by smooth bounded
functions $q_\epsilon$. Stable contraction and invariance give
\[
 \int |q(T^ny)-q(T^nx)|\,d\pi^s(x,y)\longrightarrow0.
\]
Indeed the replacement of the two occurrences of $q$ costs at most
$2C_P\|q-q_\epsilon\|_1$, uniformly in $n$, and the difference of
$q_\epsilon$ tends to zero along the contracting leaves. Take $n$ to
infinity before removing the approximation. Along these leaves the
lattice itinerary is identical, so iteration of
\eqref{eq:v35-general-phase} and the action identity give, almost
everywhere on stable pairs,
\[
 q(y)/q(x)=\exp\left(ib\int_{[x,y]^s}\vartheta\right).
\]
Backward iteration gives the same formula on unstable pairs.
The pair disintegration and Fubini's theorem permit their multiplication
around almost every small stable--unstable rectangle $Q$ in $P$.
Consequently $\exp(ib\int_Qd\vartheta)=1$.
Since $d\vartheta=\cos\varphi\,d\varphi\wedge dr$ is nonvanishing
on this rectangle, nondegenerate rectangles of arbitrarily small
positive area contradict this equality unless $b=0$. This is the
argument of Section~\ref{sec:measurable-phase-rigidity}, using the general
one-form and not a circular parametrization.

For $b=0$ both leaf ratios are one. Product absolute continuity and
Fubini imply $q=c$ almost everywhere on $P$ for some $|c|=1$.
Discarding the invariant null set on which an iterate of the phase
equation fails, the orbit saturation of $P$ has full measure. Hence
$q/c$ takes values in the countable group
\[
 \{\exp(i(u\cdot k+s n)):(k,n)\in\Z^2\times\Z\}.
\]
Define its kernel
$L=\{(k,n):u\cdot k+s n\in2\pi\Z\}$.
Enumeration of the countable quotient defines a measurable function
$h:M\to\Z^3/L$ by $q/c=\exp(i(u,s)\cdot h)$. Uniqueness in the
quotient gives the exact measurable relation
\begin{equation}\label{eq:v35-space-count-lift}
 h\circ T=h+[(\kappa,1)].
\end{equation}

If $L$ were proper, Lemma~\ref{lem:v35-finite-character} would give
$\chi(k,n)=a\cdot k+c_0 n\pmod p$, with $(a,c_0)\ne0$.
It descends to $\Z^3/L$. On the actual $p$-lattice cover, written as
$(x,l)\in M\times(\Z/p\Z)^2$, the collision map is
$\widetilde T(x,l)=(Tx,l+\kappa(x))$. Put
\[
 H(x,l)=\exp\left(\frac{2\pi i}{p}
                         (a\cdot l-\chi(h(x)))\right).
\]
Equation~\eqref{eq:v35-space-count-lift} yields
$H\circ\widetilde T=e^{-2\pi i c_0/p}H$.
If $c_0\ne0$, this is a nontrivial unit-modulus eigenvalue, excluded
by mixing of the finite cover. If $c_0=0$, then $a\ne0$ and the
sheet average of $H$ is zero, so this is a nonconstant invariant
function, excluded by ergodicity. Both cases are impossible. Thus
$L=\Z^3$, which gives the asserted values of $u,s$. Finally
\eqref{eq:v35-general-phase} reduces to invariance of $q$.
\end{proof}

\begin{corollary}[Nondegeneracy from the same arithmetic]
\label{cor:v35-general-covariance}
Suppose the bounded collision observable
$f^{\rm c}=(\kappa_1,\kappa_2,\tau-\bar\tau)$ has zero mean and
exponentially summable covariances. Its Green--Kubo matrix $\Sigma$ is
positive definite. For a compact parameter family with common
correlation bounds and continuous finite-lag covariances, the matrices
are continuous and have a common positive lower eigenvalue.
\end{corollary}
\begin{proof}
Let $v\cdot\Sigma v=0$ and put $g=v\cdot f^{\rm c}$. Exponential
summability gives $\sum_{k\ge1}k|\int g(g\circ T^k)d\nu|<\infty$.
The finite covariance expansion then bounds $\|S_mg\|_2$ uniformly
in $m$. The averages $b_M=M^{-1}\sum_{m=1}^M S_mg$ are bounded in
$L^2$ and satisfy
$(I-U)b_M=g-M^{-1}\sum_{m=1}^M U^m g\to g$, where $Ug=g\circ T$.
A weakly convergent subsequence supplies $g=b-b\circ T$ with a real
$b\in L^2$. For every real $t$, exponentiating $-itb$ gives
\eqref{eq:v35-general-phase} with frequencies
$(tv_1,tv_2)$, $-tv_3\bar\tau$, and $tv_3$.
Theorem~\ref{thm:v35-rotation-free-rigidity} forces $v_3=0$ and
$t(v_1,v_2)\in(2\pi\Z)^2$ for every $t$, hence $v=0$.
Uniform covariance convergence and finite-lag continuity give continuity
of $\Sigma$. Compactness supplies the positive minimum eigenvalue.
No common measurable representative of the transfer functions, or
uniform size of the local product rectangle, is needed for this last
compactness argument.
\end{proof}

\begin{remark}[The section occupation is not inserted into the twist]
\label{rem:v35-return-distinction}
The integer collision step in \eqref{eq:v35-space-count-lift} is the
constant $1$, not the discontinuous occupation $\one_{Y^*}$.
The theorem strengthens the arithmetic used for the physical
collision--roof spectrum. It does not provide a Lasota--Yorke estimate
for the accumulated section occupation or for unbounded induced
return twists. The common pointwise return correction and the return
frequency complement in Theorem~\ref{thm:LLT} remain distinct analytic
requirements for the original four-coordinate problem.
\end{remark}
